\clearpage
\VerdierPhysicalPageAnchor{0042}%
\VerdierSemanticAnchor{ch1-1-3-7-cont}
$G^\circ$-圏の言語だけを用いる。しかし「次数付き圏」の観点を完全には忘れず、
次の用語を用いる。

\VerdierSemanticAnchor{ch1-1-3-8}
\noindent\textbf{定義 1.3.8.}\quad
$G$ を群、$(\mathcal C,G^\circ,T,c)$ を $G^\circ$-圏とする。$g,h$ を $G$ の
二元、$X,Y$ を $\mathcal C$ の二対象とする。$\mathcal C$ における射

\VerdierSemanticAnchor{ch1-1-3-8-morphism}
\[
  u\colon T(h)X\longrightarrow T(gh)Y
\]
を、\emph{$X$ から $Y$ への次数 $g$ の射}という。

この用語が混乱を生じないことを示そう。実際、$G$ の各元対 $h,h'$ に対し、
関手 $T(h^{-1}h')$ と同型 $c(h,h^{-1}h')$, $c(gh,h^{-1}h')$ は双関手的同型

\VerdierSemanticAnchor{ch1-1-3-8-eq1}
\begin{equation}
  i_g(h,h')\colon
  \operatorname{Hom}_{\mathcal C}\bigl(T(h)X,T(gh)Y\bigr)
  \overset{\sim}{\longrightarrow}
  \operatorname{Hom}_{\mathcal C}\bigl(T(h')X,T(gh')Y\bigr)
  \tag{1.3.8.1}
\end{equation}
を定める。

余サイクル恒等式 (1.2.1.3) により、各 $h,h',h''\in G$ に対して

\VerdierSemanticAnchor{ch1-1-3-8-eq2}
\begin{equation}
  i_g(h,h'')=i_g(h',h'')i_g(h,h')
  \tag{1.3.8.2}
\end{equation}
となる。したがって、同型 $i_g(h,h')$ によって種々の集合

\VerdierSemanticAnchor{ch1-1-3-8-identification}
\[
  \operatorname{Hom}_{\mathcal C}\bigl(T(h)X,T(gh)Y\bigr)
\]
を同一視すれば、式 (1.3.8.2) はこれらの同一視が両立することを示す。さらに、
任意の $g,g',h,h'\in G$ と $\mathcal C$ の任意の三対象 $X,Y,Z$ に対し、
次の図式 (1.3.8.3) は可換である：

\VerdierSemanticAnchor{ch1-1-3-8-eq3}
\begin{center}
  \resizebox{0.92\textwidth}{!}{%
    $\begin{tikzcd}[ampersand replacement=\&,column sep=5em,row sep=4em]
      \operatorname{Hom}_{\mathcal C}\bigl(T(h)X,T(gh)Y\bigr)
      \times
      \operatorname{Hom}_{\mathcal C}\bigl(T(gh)Y,T(g'gh)Z\bigr)
      \arrow[r]
      \arrow[d,
        "{i_g(h,h')\times i_{g'}(gh,gh')}"',
        "\sim"]
      \&
      \operatorname{Hom}_{\mathcal C}\bigl(T(h)X,T(g'gh)Z\bigr)
      \arrow[d,"{i_{g'g}(h,h')}"',"\sim"]
      \\
      \operatorname{Hom}_{\mathcal C}\bigl(T(h')X,T(gh')Y\bigr)
      \times
      \operatorname{Hom}_{\mathcal C}\bigl(T(gh')Y,T(g'gh')Z\bigr)
      \arrow[r]
      \&
      \operatorname{Hom}_{\mathcal C}\bigl(T(h')X,T(g'gh')Z\bigr)
    \end{tikzcd}$%
  }
  \hfill\textup{(1.3.8.3)}
\end{center}

\VerdierSemanticAnchor{ch1-1-3-8-composition}
（横の矢印は $\mathcal C$ における合成で定義される。）したがって図式
(1.3.8.3) により、次数がそれぞれ $g,g'$ の二射の合成を曖昧さなく、より正確
には同型 $i_g(h,h')$ による同一視のもとで曖昧さなく定義できる。この合成の
定義が
